% Figure: axisymmetric thermal stress in a solid disk with a hotter
% centre and cooler rim (a parabolic temperature drop T(r)=T0(1-(r/a)^2)).
% The radial stress is compressive over the interior and zero at the free
% rim; the hoop stress is compressive near the centre and swings tensile at
% the rim, which is the crack-driving stress. Both are normalised by
% E alpha T0 and plotted against r/a. Closed-form for this T-profile:
%   sigma_r/(E alpha T0)    = (1/8)( (r/a)^2 - 1 )     [<=0, zero at rim]
%   sigma_theta/(E alpha T0)= (1/8)( 3 (r/a)^2 - 1 )   [tensile for r/a>1/sqrt3]
\begin{figure}[htbp]
\centering
\begin{tikzpicture}
\begin{axis}[
  width=11cm, height=6.5cm,
  xlabel={normalised radius $\ r/a$},
  ylabel={stress / $(E\alpha T_0)$},
  xmin=0, xmax=1.0, ymin=-0.18, ymax=0.30,
  axis lines=left,
  grid=major, grid style={enggray!20},
  tick label style={font=\scriptsize},
  label style={font=\small},
  legend style={font=\scriptsize, at={(0.02,0.98)}, anchor=north west, draw=enggray!40},
]
% radial stress: (1/8)(x^2 - 1), compressive, zero at rim
\addplot[engblue, very thick, domain=0:1, samples=80]
  {(1/8)*(x*x - 1)};
\addlegendentry{radial $\sigma_r$}
% hoop stress: (1/8)(3x^2 - 1), tensile near rim
\addplot[engred, very thick, domain=0:1, samples=80]
  {(1/8)*(3*x*x - 1)};
\addlegendentry{hoop $\sigma_\theta$}
% zero line
\addplot[enggray, dashed, domain=0:1, samples=2] {0};
% mark the rim tensile hoop value (1/4 at x=1)
\addplot[engamber, only marks, mark=*, mark size=1.6pt] coordinates {(1.0,0.25)};
\node[engamber, font=\scriptsize, anchor=east] at (axis cs:0.97,0.255)
  {rim hoop tension $\tfrac{1}{4}E\alpha T_0$};
% mark sign change of hoop at r/a = 1/sqrt3
\node[engred, font=\scriptsize, anchor=south west] at (axis cs:0.58,0.01)
  {hoop turns tensile};
\end{axis}
\end{tikzpicture}
\caption[Thermal stress in a heated disk]{Axisymmetric thermal stress in a
solid disk whose centre is hotter than its rim, for the parabolic
temperature drop $T(r)=T_0\,[1-(r/a)^2]$. The stresses are normalised by
$E\alpha T_0$ and plotted against the normalised radius. The radial stress
(blue) is compressive everywhere inside and falls to zero at the free rim,
as a free edge requires. The hoop stress (red) is compressive near the hot
centre, which wants to expand and is held back by the cooler ring around it,
and swings into tension toward the rim, reaching $\tfrac{1}{4}E\alpha T_0$ at
the edge (amber marker). That rim tension is the crack-driving stress: a
disk heated rapidly at its centre, or a brake rotor or a grinding wheel run
hot, cracks radially from the rim inward, not from the hot centre. No
external load is applied anywhere; the entire stress field is the
self-equilibrating consequence of a non-uniform temperature in a body that
is not free to expand differentially.}
\label{fig:vol03ch09:thermal-disk}
\end{figure}
